Softwaresysteme neigen dazu, im Laufe ihrer Entwicklung zu wachsen; wie sich dieses Wachstum in der Basis einer langlebigen Linux-Distribution äussert, ist jedoch wenig untersucht. Diese Arbeit liefert eine reproduzierbare, releaseweise Charakterisierung des persistenten Kerns der minbase-Installation von Debian, der minimalen Paketmenge, die \texttt{debootstrap} für jedes Release auflöst: der 25 Quellpakete, die in mindestens 11 der 13 stabilen Releases von Potato (2000) bis Trixie (2025) in dieser Menge enthalten sind. Eine containerisierte Pipeline lädt die Paketindizes und Quellarchive jedes Releases herunter und berechnet fünf Metriken pro Release: Installationsgrösse, Abhängigkeitsstruktur, Quelltextzeilen, zyklomatische Komplexität pro Funktion und das Volumen distributionsspezifischer Patches. Die Ergebnisse beschreiben ein Wachstum durch Anlagerung. Der persistente Kern wächst etwa um das 5,3-Fache in Quelltextzeilen und etwa um das Vierfache in der Installationsgrösse, während der Median der zyklomatischen Komplexität pro Funktion in allen Releases bei 3 verbleibt und der Median der Abhängigkeiten pro Paket in jedem Release nach Potato bei 1 liegt; die Zusammensetzung der minbase-Installation ändert sich vor allem durch die Paketierungsstruktur und nicht durch den Funktionsumfang. Die Patch-Metrik erweist sich erst ab dem Quellformat 3.0 (quilt) als vergleichbar und nimmt innerhalb dieses Fensters ab. Die Studie ist explorativ-deskriptiv angelegt und behauptet keine kausalen oder korrelativen Zusammenhänge zwischen den Dimensionen. Die Pipeline und eine interaktive Begleitanwendung werden zusammen mit den Ergebnissen veröffentlicht, um Reproduktion und Erweiterung zu ermöglichen.

\vspace{6em} 

\noindent\textbf{Schlüsselwörter:} Debian, Softwareevolution, Software-Bloat, Softwaremetriken, zyklomatische Komplexität, Mining Software Repositories
